% Owner: Khalid  |  Task: T15  |  Drafted by ORCA agents, finalised by the owner.
% Rules: prose only; every claim cited with \citet/\citep to a literature/sources.yaml record id;
% compare methods (G6); strengths/weaknesses/applicability (G8). See tasks/todo.md.
\section{Conclusion}
\label{sec:conclusion}

This survey reviewed knowledge distillation for image classification from 2022 to 2026. Its method set was
not chosen by hand. Three searches and a gap pass over the main vision and machine-learning venues, two
journals and recent arXiv preprints produced 112 methods that meet fixed inclusion criteria, and the
exclusion of every assessed paper is recorded with a code.

The methods fall into five families. Logit methods are the largest group and the cheapest to run: they need
no projector and apply to any pair with a shared label space. Feature and relation methods give some of the
highest single-method CIFAR-100 numbers, but they need access to intermediate layers and a module per pair,
and the best of them can enlarge the deployed student \citep{chen2022simkd}. Training schemes trade accuracy
on the standard pairs for removing or reducing the cost of a fixed pre-trained teacher. Reliability methods
decide which teacher outputs to trust, or retrain the teacher for the student. A smaller group changes how the
task and distillation gradients are combined \citep{zhao2023dot,huang2025deepkd,hayder2026dtokd}. We
reviewed it next to the general multi-objective rules it borrows from, and placed priority-based rules such as
PCD \citep{varam2026pcd} among them as one option, without reported numbers on the standard pairs.

The compiled results support a cautious reading. On CIFAR-100 ResNet32$\times$4$\to$ResNet8$\times$4 the best
reported configurations are about six points above classical KD, yet most single methods lie within two
points of one another, and on ImageNet nearly all rows sit within two points. Many of the highest rows stack
one method on another, change the teacher, or add data, and the papers differ in how many runs they average.
Gaps of a few tenths of a point therefore say little. A second reading is that optimizer-side changes stack on
loss-side ones, so they should be tested on top of both plain KD and a strong logit loss.

Several questions stay open. Cross-paper comparison would improve if papers reported seed statistics and
stand-alone numbers for stacked methods. Results depend on the training recipe, as patient and consistent
teaching shows \citep{beyer2022patientkd}, and a student can generalise well while matching the teacher
poorly \citep{stanton2021kdwork}. For the objective-combination methods, the main untested point is how
often the task and distillation gradients conflict on the standard benchmarks at all, and whether the
extra gradient computations pay off when they do. The debate on whether a tuned weighted sum is enough in
multi-task learning \citep{xin2022mtohelp,kurin2022unitary,royer2023scalarization} suggests that strong
weighted-sum baselines belong in every such comparison. Section~\ref{sec:future} discusses these
directions in more detail.
